%% ============================================================
%%  Section: Introduction  (2026-06-12, post-JMCB revision)
%%  Depositor Discipline in Community Banking
%%
%%  Structure:
%%    ¶1   Hook + explicit research question
%%    ¶2   Identification challenges + CECL solution + 2023-stress-by-design preview
%%    ¶3   Main results: within-bank composition shift (core), level effect, magnitudes
%%    ¶4   Where the deposits go: insured wholesale substitution + local spillovers
%%    ¶5   Signal validation (large banks) — demoted to suggestive
%%    ¶6   Financial consequences + lending + sophistication, with magnitudes
%%    ¶7   Contribution 1: depositor discipline / market discipline lit
%%    ¶8   Contribution 2: accounting disclosure + CECL consequences lit
%%    ¶9   Contribution 3: deposit insurance design + policy
%%
%%  Changes from JMCB version (per revision-plan-20260612.md):
%%    - removed "absent a crisis trigger" and "non-crisis setting" framing
%%    - identification ¶ previews flexible controls + cohort-timing tests (RA2/RB1)
%%    - composition shift promoted to headline result (survives all controls)
%%    - signal-validation ¶ compressed; 191 bps point estimate dropped (RA3)
%%    - new ¶4 on deposit destination and spillovers (RB2)
%%    - lending result added (RA8); magnitudes sentences added throughout
%%    - "predetermined" language softened (RA1)
%% ============================================================

\section{Introduction}
\label{sec:intro}

Whether uninsured depositors discipline bank risk remains one of the central open
questions in banking.  Theory predicts they should: uninsured depositors bear
full downside risk and have strong incentives to withdraw or demand higher rates
when they detect deterioration in bank financial condition
\citep{calomiris1991role, diamond2001liquidity}.  A large empirical literature
documents that riskier banks pay higher deposit rates and experience larger
uninsured outflows \citep{park1995market, martinez2001depositors,
maechler2006dynamic, iyer2016tale, egan2017deposit, chen2024liquidity}, yet whether these patterns
reflect a causal disciplinary response or equilibrium correlation remains
unsettled.  The question is pressing: against a backdrop of expanding government
protection for uninsured depositors, the Main Street Depositor Protection Act
(S.~2999, reintroduced March 2026) would formally extend coverage to
\$10~million per account, making it important to understand whether this
discipline operates before the incentives that sustain it are further weakened. Historical
evidence suggests the stakes are large, as early state deposit insurance programs in the
United States eliminated market discipline, enabling insured banks to take on excessive
risk \citep{calomiris2019stealing}.
We ask whether uninsured depositors respond to mandatory credit-risk disclosures
by withdrawing balances or demanding higher rates, and we provide causal evidence
on this question using U.S.\ community banks, where institutions below
\$10~billion in assets carry no credible government backstop and identification
conditions are favorable \citep{o1990deposit, flannery1996evidence}.

Achieving clean identification requires resolving three problems that prior work
leaves open.  First, simultaneity: a bank's risk profile, deposit composition,
and depositor behavior co-evolve, so regressions of deposit flows on risk
measures capture equilibrium sorting rather than isolated disciplinary responses.
Second, accounting opacity: under the incurred-loss model, banks could delay
credit loss recognition, systematically understating risk in public disclosures
and limiting depositors' ability to monitor \citep{bushman2012accounting,
bushman2015delayed}.  Third, confounding from aggregate stress: the periods with
the greatest variation in bank distress are precisely those in which government
interventions and depositors' own liquidity needs contaminate the measured
response \citep{bennett2015market}.  We address the first two problems by
exploiting mandatory CECL adoption for community banks in 2023 as an information
shock whose cross-sectional intensity was largely fixed before the event.  CECL
requires banks to recognize lifetime expected credit losses on existing loan
portfolios at adoption, producing a one-time adjustment to retained earnings
whose magnitude reflects loan vintages assembled years earlier.  Banks knew the
standard was coming, but the size of the Day-One charge---the treatment
variation we use---was pinned down by the composition and seasoning of the
existing loan book rather than by contemporaneous deposit-market conditions.
The third problem we confront directly rather than assume away: community banks
adopted CECL in the same quarter that the failures of Silicon Valley Bank and
Signature Bank put bank risk on every depositor's mind.  We show that our
estimates are robust to allowing deposit dynamics to vary flexibly with the full
set of 2022Q4 bank characteristics that governed exposure to the 2023 stress---%
uninsured deposit shares, size, capital, asset quality, profitability,
commercial-real-estate concentration, and unrealized securities losses---each
interacted with calendar-quarter fixed effects, and that the response patterns
line up with each cohort's own adoption quarter rather than with March 2023.
\citet{kim2026current} and \citet{gee2022decision}
confirm that the Day-One adjustment is informative, predicting subsequent
charge-offs and improving analyst forecast accuracy; \citet{nier2006market} and
\citet{chen2022bank} establish that disclosure improvements of this type enable
depositors to act on credit risk signals.

Our central result, estimated on a sample of more than 4,000 community banks over
quarterly observations from 2016 to 2025, is a within-bank shift in the
composition of time deposits away from uninsured and toward insured balances.  A
triple-difference specification stacks the insured and uninsured series in a
single panel and adds bank-by-quarter fixed effects, absorbing all bank-specific
time-varying shocks---including any bank-level effect of the 2023 stress; the
triple interaction implies that each percentage point of CECL exposure reduces
uninsured time deposits by $0.18$ percentage points of assets \emph{relative to}
insured time deposits within the same bank and quarter.  This composition result
is the most robust fact in the paper: it survives the characteristics-by-quarter
controls, a fully saturated specification, alternative treatment scalings and
cutoffs, and the exclusion of the largest-adjustment bank.  The simple
difference-in-differences levels are consistent with it.  Banks in the right
tail of CECL exposure (Day-One adjustment of at least 2.5 percent of
pre-adoption equity) experienced a 0.394 percentage-point contraction in
uninsured time deposits relative to total assets---roughly 7 percent of the
sample mean, or about \$1.3~million for the median treated bank---while insured
time deposits rose by 0.71 percentage points.  Event studies confirm that
deposit trajectories were parallel across treatment and control banks before
adoption and that the divergence emerged at adoption and persisted through the
end of the sample.  We are transparent that the \emph{level} of the uninsured
decline attenuates when deposit-composition characteristics are allowed their
own time paths---as it mechanically must, since those controls absorb part of
the uninsured channel itself---which is why we anchor the paper on the
within-bank composition contrast that no such control can explain away.

We also answer the question any composition result raises: where do the deposits
go?  High-CECL banks replace the lost uninsured balances with \emph{insured}
wholesale funding---insured brokered deposits rise by roughly 0.5 percentage
points of assets and reciprocal deposits, which split large accounts across
network banks to obtain full coverage, also increase---while uninsured brokered
funding does not move.  At the same time, branch-level data from the FDIC
Summary of Deposits show that deposit growth after 2023 is faster at low-CECL
banks operating in counties more exposed to high-CECL competitors, direct
evidence that part of the withdrawn funding re-intermediates locally.  Both
margins are exactly what a disciplinary mechanism predicts: depositors and the
banks competing for them reprice and reallocate around the newly disclosed
credit risk, partly by purchasing insurance through market arrangements
\citep{kim2024insurance}.

Before turning to deposits, we verify that the Day-One disclosure moves
informed-market prices at all, using the cohort of large, publicly traded banks
that adopted CECL in 2020Q1---the only cohort with observable equity prices and
credit default swap spreads around adoption.  Market capitalization falls and
CDS spreads widen differentially for high-CECL large banks beginning in the
adoption month, with flat pre-adoption leads.  Because that cohort adopted on
the eve of the pandemic, we treat this evidence as suggestive corroboration
that the adjustment carries credit-relevant news, not as a standalone causal
estimate.

The deposit reallocation carries real financial costs.  High-CECL banks paid
approximately 8--9 basis points more annually in interest expense relative to
assets---about \$0.3~million per year for the median treated bank---consistent
with paying up to replace outflowing uninsured balances with insured wholesale
funding.  Return on assets fell by roughly 15 basis points per year and net
interest margin compressed by approximately 9 basis points annually, confirming
that higher funding costs were not passed through to borrowers.  The pressure
extends to the asset side: quarterly loan growth at high-CECL banks slows by
0.65 percentage points, concentrated in commercial real estate, linking the
funding shock to credit supply in the spirit of \citet{granja2025current}.  The
effects are amplified by depositor sophistication.  Splitting the sample by a
deposit-weighted branch-market index of financial sophistication, constructed
from FDIC branch data merged to ZIP-code American Community Survey estimates of
educational attainment and IRS investment-income filing rates, we find that
high-CECL banks in high-sophistication markets paid an additional 17--18 basis
points annually in interest expense and suffered an additional 30 basis point
drag on return on assets; the pattern is unchanged when the index is built from
pre-period (2019) branch weights, addressing the concern that branch footprints
responded to treatment.

This paper contributes first to the empirical literature on depositor and market
discipline in banking.  A substantial body of evidence documents that riskier
banks pay higher rates and experience larger uninsured outflows
\citep{park1995market, martinez2001depositors, maechler2006dynamic,
egan2017deposit}, and that market-priced signals contain information about bank
distress \citep{flannery1996evidence, nier2006market}.  The causal
interpretation of these patterns has proven difficult to establish, however.
\citet{bennett2015market} and \citet{berger2015depositors} document that
government guarantees attenuate estimated effects in the crisis episodes that
generate the most variation, and institutions subject to implicit government
support carry systematically smaller estimated discipline effects
\citep{jacewitz2018deposit}.  Recent work on the 2023 episode shows that
depositors ran on banks with high uninsured shares and unrealized losses
\citep{caglio2024depositors}; our contribution is orthogonal variation
\emph{within} that high-attention period: a bank-specific, predetermined
disclosure whose intensity is essentially uncorrelated with the balance-sheet
characteristics that governed run exposure, so that the disciplinary response
to new credit-risk information can be separated from the aggregate flight to
safety.

This paper also contributes to the accounting and disclosure literature on CECL
and its economic consequences.  \citet{bushman2012accounting} and
\citet{bushman2015delayed} establish that delayed loss recognition under the
incurred-loss model reduced bank transparency and attenuated market discipline;
our findings provide direct evidence that the mandated switch to forward-looking
expected-loss accounting triggered the disciplinary response that opacity had
been suppressing.  Prior work on CECL's real effects focuses on lending behavior
\citep{granja2025current}, provisioning dynamics \citep{kim2026current}, and
decision-usefulness \citep{gee2022decision}.  Unlike these studies, which examine
bank-side responses or analyst outcomes, we trace the full causal chain from
mandatory accounting disclosure to depositor behavior to bank funding costs,
profitability, and lending, documenting the funding-market channel through which
improved disclosure standards alter the cost of risk-taking.

Finally, this paper speaks to the debate over deposit insurance design and the
scope of government guarantees.  Prior work establishes that broader insurance
coverage reduces the incentive of uninsured depositors to monitor and exert
pressure on bank risk-taking, weakening the market discipline that complements
regulatory oversight \citep{demirgucc2004market, flannery2019market}.  Our
evidence that uninsured depositors monitor and respond to credit risk
information at community banks supports the view that depositor discipline is a
functioning governance mechanism in the institutional setting where theory
predicts it should work best---and the substitution into insured brokered and
reciprocal deposits shows that much of the response takes the form of
repurchasing insurance rather than exiting the banking system
\citep{kim2024insurance}.  Policies that move a substantial volume of currently
uninsured deposits inside the insurance perimeter should be evaluated not only
for their run-prevention benefits but also for the monitoring and pricing
discipline they displace.
